\section{Introduction}
Anomaly detection (AD) aims to identify rare or abnormal instances that deviate from an underlying notion of "normality" \citep{hawkins1980identification, chandola2009anomaly}. This problem is central in high-stakes applications where anomalies may indicate fraudulent transactions, cybersecurity intrusions, safety issues, clinical outliers, or abnormal events in surveillance video \citep{ngai2011application, buczak2015survey, vu2019robust, fernando2021deep, landauer2023deep}. In practice, data in these domains are often tabular and heterogeneous: each record may mix numerical measurements, categorical attributes, and increasingly, free-form text with missing values and non-trivial cross-field dependencies. These characteristics make unsupervised detection difficult, as standard tabular AD pipelines typically assume fully structured features and either handle text only crudely or rely on heavy feature engineering that can remove useful semantics and miss complex interactions. Large language models (LLMs) have recently become an attractive foundation for mixed-type tabular AD because they can operate on a serialized representation of each row and naturally exploit the semantics of categorical values and raw text \citep{dinh2022lift, hegselmann2023tabllm, fang2024large}. Viewed as a serialized sequence, each row can be processed by an LLM in a unified way, allowing the model to exploit the semantics of categorical values and raw text with minimal preprocessing. However, applying LLMs to unsupervised tabular AD is non-trivial. Concretely, the serialization should respect basic tabular invariances (e.g., column order), numerical fields require representations that are robust to tokenizer-induced distortions, and likelihood-based scores should be calibrated to avoid systematic bias from variable-length text.

A foundational early work that first demonstrated the practicality of applying LLMs to tabular anomaly detection is AnoLLM \citep{tsai2025anollm}. AnoLLM serializes each row into a text template, discretizes numerical values into bins, and fine-tunes a LLM on normal data via next-token prediction. At inference time, it scores anomalies using negative log-likelihood, together with procedures to better handle mixed-type tables. Despite this progress, AnoLLM can be viewed as a one-shot SFT-style adaptation of a generative model to normal data, leaving a meaningful gap for unsupervised anomaly detection. Maximizing likelihood on normal training data provides only a weak signal for tightening the boundary of normality: once model fits the dominant regularities of the training corpus, further gains in separating subtle, near-normal anomalies are often limited and performance can plateau. These limitations motivate an iterative refinement mechanism that more directly reduces the mismatch between model-induced distribution and the true normal-data distribution, while remaining compatible with likelihood-based detection.

Self-play provides a principled route to such iterative refinement. Recent advances such as SPIN \citep{chen2024self} improve an SFT-initialized model by repeatedly contrasting real data with the model's own generations: at each iteration, the current model produces synthetic samples, and the next update is driven by a discriminator-like objective that favors genuine data over model-generated counterparts, providing a learning signal without additional annotations. In unsupervised tabular AD, this provides an iterative refinement signal: a critic distinguishes real normal samples from model-generated ones, and the policy is updated so that future generations better match the normal-data distribution. Unlike SPIN, which optimizes a single fixed IPM-like divergence, we extend this idea to a broader family of $f$-divergences. Different divergences induce different distribution-matching behavior---for example, KL is more mode-covering, whereas reverse KL is more mode-seeking---which is useful in anomaly detection because the geometry of the normal data manifold can vary substantially across domains. Repeating this loop can improve the quality and calibration of likelihood-based anomaly scores without requiring anomaly labels. 

In this paper, we introduce \textbf{DiSPaT}, a \textbf{Di}vergence-driven \textbf{S}elf-\textbf{P}l\textbf{a}y framework for unsupervised \textbf{T}abular anomaly detection with LLMs. \textbf{DiSPaT} treats each mixed-type row as a serialized token sequence and iteratively refines the model's understanding of normality through an adversarial self-play game. Formally, we optimize the LLM to minimize the $f$-divergence between the empirical distribution of normal data and the model-induced distribution. Conceptually, this is achieved via an alternating optimization loop: in the \textit{critic step}, a critic contrasts real normal sequences with samples from the current policy; in the \textit{policy step}, the LLM is updated using this contrastive signal so that its induced distribution moves closer to the normal-data distribution. In the final implementable objective, the critic is absorbed into an implicit log-ratio parameterization, yielding a single optimization step per iteration. Through this iterative process, the model progressively aligns its induced distribution with the true normal-data distribution. Finally, anomalies are identified based on the likelihood scores of the refined model. Our contributions are summarized as follows:



\begin{itemize}
    \item We propose a principled theoretical framework for unsupervised tabular anomaly detection via LLMs by reformulating the training objective as minimizing $f$-divergence between normal data distribution and model's generations. This formulation transcends the limitations of one-shot SFT, providing a mechanism that explicitly targets the distributional mismatch between model and true data manifold.

    \item We introduce \textbf{DiSPaT}, a practical algorithm that implements this framework via a self-play loop with an implicit critic. It preserves the critic--policy interpretation while reducing each iteration to a single tractable optimization step, enabling the model to capture complex tabular dependencies without anomalous supervision.

    \item Extensive evaluations on mixed-type benchmarks demonstrate that the proposed approach outperforms state-of-the-art baselines, including recent LLM-based and specialized deep anomaly detectors. Furthermore, detailed ablation studies on iteration depth and divergence choices validate the efficacy of the self-play mechanism.

\end{itemize}
